%% P09 --- Wartości własne, formy kwadratowe i dodatnia określoność
%%
%% Zalążek. Poniższy opis jest kontraktem tego programu: co ma uzasadnić i co
%% ma dać czytelnikowi. Napisanie programu oznacza usunięcie bloku
%% \programstub{} i zastąpienie go programem. Jeśli po przerobieniu materiału
%% opis okaże się błędny, popraw go i odnotuj sprzeczność w CLAUDE.md w sekcji
%% *Resolved questions*.

\program{Wartości własne, formy kwadratowe i dodatnia określoność}
\label{prog:P09}

\programstub{%
\textit{[Opis do przetłumaczenia --- patrz wydanie angielskie.]}

Argument: an eigenvector is a direction the matrix only stretches; a
symmetric matrix has a full orthogonal set of them (the spectral theorem,
stated and used, not proved); a quadratic form is a bowl, a saddle or a
ridge, and its eigenvalues say which. Payoff: a covariance matrix is
symmetric and positive semi-definite, so PCA is an eigen-decomposition and
not a black box; the spectral norm as the largest stretch, which is the
Lipschitz constant that bounds how much a layer can amplify; the Hessian's
eigenvalues are the shape of the loss basin, which is where \enquote{ravine}, \enquote{sharp minimum\enquote{ and }the learning rate is too high} all become one statement ---
collected here and spent in P16 and P19.

\medskip
\textbf{Planowana długość:} 65 ramek.
}
